% -===============                                    ===============%
% -===============                                    ===============%
% -===============                                    ===============%
%                           PACKAGES                                 %
% -===============                                    ===============%
% -===============                                    ===============%
% -===============                                    ===============%
% shut up, warnings
\usepackage{sectsty}
\let\underbar\relax
\usepackage{silence}
\WarningFilter*{hyperref}{Token not allowed in a PDF string}
\WarningsOff[caption]
\usepackage[utf8]{inputenc}
\usepackage[T1]{fontenc}
\usepackage{tabularx}
\usepackage{hyperref}
\usepackage{makecell}
\usepackage{lipsum}
\usepackage{algorithm}
\usepackage{minted}
\usepackage{algpseudocode}
\usepackage{upquote}
%\usepackage[titlenumbered,ruled,noend,algo2e]{algorithm2e}
\usepackage{url}
\usepackage{placeins}
\usepackage{booktabs}
\usepackage{soul} % strikethrough, underline, highlight
\usepackage{amsfonts}
\usepackage{nicefrac}
\usepackage{microtype}
%\usepackage[numbers]{natbib}
\usepackage{proba} % \P, \E, \Var, etc.
\usepackage{wrapfig}
\usepackage{caption}
\usepackage{amssymb}
\usepackage{pifont}
\newcommand{\cmark}{\ding{51}}
\newcommand{\xmark}{\ding{55}}
\usepackage{subcaption}
\usepackage{float}
\usepackage{rotating}
\usepackage{etoolbox}
\usepackage{xparse}
\usepackage{grffile}
\usepackage{amsmath}
\usepackage{xspace}
\usepackage{euscript}
\usepackage{enumitem}
\usepackage{mathtools}
\usepackage{tikz}
\usepackage{tablefootnote}
\usepackage[flushleft]{threeparttable}
\usepackage{multirow}
\usepackage[title]{appendix}
\usepackage{arydshln}
\usepackage{array, booktabs}
\usepackage{comment}
\usepackage{graphicx}
\usepackage{adjustbox}
\usepackage{dsfont} % \mathds for blackboard bold
\usepackage{amsmath,amsthm}
\usepackage{balance}
\usepackage{multicol}
\usepackage{xcolor}
\usepackage{nameref}
\usepackage{pdfpages}
\usepackage{braket}
\usepackage[normalem]{ulem} % normalem so \emph stays italic
\usepackage{bbding}
\usepackage[capitalise,noabbrev]{cleveref}
\usepackage{xfrac}
\usepackage[usestackEOL]{stackengine}
\usepackage{indentfirst}
\usepackage{csquotes}
\usepackage{pgf}
\usepackage{varwidth}
\usepackage{verbatimbox}
\usepackage{verbatim}
\usepackage{bm}
\usepackage{pythonhighlight}
\usepackage[listings]{tcolorbox} % fancy boxes
\tcbuselibrary{breakable, listings, skins}
\usepackage[framemethod=default]{mdframed} % simpler boxes

% -===============                                    ===============%
% -===============                                    ===============%
% -===============                                    ===============%
%                           CAPTION SETTINGS                         %
% -===============                                    ===============%
% -===============                                    ===============%
% -===============                                    ===============%

% keep captions readable, not giant
\captionsetup{font=small}
\captionsetup[sub]{font=small}
\captionsetup[subfigure]{labelformat = parens, labelsep = space, font = small}
% shortcuts for subfigure references in caption text
\newcommand{\figleft}{{\em (Left)}}
\newcommand{\figcenter}{{\em (Center)}}
\newcommand{\figright}{{\em (Right)}}
\newcommand{\figtop}{{\em (Top)}}
\newcommand{\figbottom}{{\em (Bottom)}}
\newcommand{\captiona}{{\em (a)}}
\newcommand{\captionb}{{\em (b)}}
\newcommand{\captionc}{{\em (c)}}
\newcommand{\captiond}{{\em (d)}}

% -===============                                    ===============%
% -===============                                    ===============%
% -===============                                    ===============%
%                         COLOR MACROS                               %
% -===============                                    ===============%
% -===============                                    ===============%
% -===============                                    ===============%

\definecolor{pastelblue}{RGB}{76,113,175}
\definecolor{pastelgreen}{rgb}{0.8, 1.0, 0.8}
\definecolor{pastelred}{rgb}{1.0, 0.8, 0.8}
\definecolor{pastelgrey}{RGB}{230,230,230}
\definecolor{pastelbeige}{RGB}{243,236,221}
\definecolor{pastelpurple}{RGB}{154,139,192}
\definecolor{salmon}{RGB}{250, 128, 114}
\definecolor{darkgreen}{rgb}{0,0.6,0}
\definecolor{darkred}{rgb}{0.5,0,0}
\definecolor{verylightgreen}{HTML}{F6FFF9}
\definecolor{verylightred}{HTML}{FFF4F3}
\definecolor{verylightgray}{HTML}{F4F6F6}
\definecolor{navyblue}{HTML}{473992}
\definecolor{redorange}{HTML}{ED135A}
\definecolor{babyblueeyes}{rgb}{0.63, 0.79, 0.95}
\definecolor{lightpink}{rgb}{1.00, 0.714, 0.757}
\definecolor{peach}{rgb}{1.0, 0.9, 0.71}
\definecolor{pastelpink}{rgb}{1.0, 0.82, 0.86}
\definecolor{deepchampagne}{rgb}{0.98, 0.84, 0.65}

% green/red inline highlight shorthands
\newcommand{\gb}[1]{\colorbox{pastelgreen}{#1}}
\newcommand{\rb}[1]{\colorbox{pastelred}{#1}}


% -===============                                    ===============%
% -===============                                    ===============%
% -===============                                    ===============%
%                           TIKZ MACROS                              %
% -===============                                    ===============%
% -===============                                    ===============%
% -===============                                    ===============%
\usepackage{tikzpeople} % stick figures in DAGs
\usetikzlibrary{shapes,decorations,arrows,calc,arrows.meta,fit,positioning}

% default DAG style — causal graph convention
\tikzset{
    -Latex,auto,node distance =1 cm and 1 cm,semithick,
    state/.style ={ellipse, draw, minimum width = 0.7 cm},
    point/.style = {circle, draw, inner sep=0.04cm,fill,node contents={}},
    bidirected/.style={Latex-Latex,dashed},
    el/.style = {inner sep=2pt, align=left, sloped}
}

% Warning symbol
\newcommand\Warning{%
 \makebox[1.4em][c]{%
 \makebox[0pt][c]{\raisebox{.1em}{\small!}}%
 \makebox[0pt][c]{\color{red}\Large$\bigtriangleup$}}}%

%                    LATIN JARGONS & TEXT HELPERS                     %
% -===============                                    ===============%
% -===============                                    ===============%
% -===============                                    ===============%

\def\eg{\textit{e.g.}, }
\def\ie{\textit{i.e.}, }
\def\etc{\textit{et cetera}}
\def\cf{\textit{cf.}\ }

% Annotation and editing helpers
\renewcommand{\paragraph}[1]{{~\\
\noindent \textbf{#1.}}} % tighter paragraph spacing
\newcommand{\Appendix}{supplement\xspace}
\newcommand{\todo}[1]{{\color{red}[\textbf{TODO:} #1]}}
\newcommand{\leo}[1]{{\color{blue}[Leo:#1]}}
\newcommand{\citehere}[0]{{\color{red}
$^{\text{cite}}$}} % "will cite later" marker

% Text box

\newcommand{\centertikztextbox}[2]{{
\begin{tikzpicture}
\node[draw, align=center,fill=#1,thick](3) at (0,0) {
\begin{varwidth}{0.9\linewidth}
#2
\end{varwidth}
};
\end{tikzpicture}
}
}

\newcommand{\tikztextbox}[2]{{
\begin{flushleft}
\begin{tikzpicture}
\node[draw, align=left,fill=#1,thick](3) at (0,0) {
\begin{varwidth}{0.9\linewidth}
#2
\end{varwidth}
};
\end{tikzpicture}
\end{flushleft}
}
}
% Dashed version
\newcommand{\dashedtikztextbox}[3]{{
\begin{flushleft}
\begin{tikzpicture}
\hspace{#2}
\node[draw, align=left,fill=#1, thick, dashed](3) at (0,0) {
\begin{varwidth}{0.9\linewidth}
#3
\end{varwidth}
};
\end{tikzpicture}
\end{flushleft}
}
}

% Arrow: [ style arguments, eat up, eat down, move right ]
\newcommand{\tikzarrow}[4]{{
\vspace{-#2}
\begin{flushleft}
\begin{tikzpicture}
\hspace{#4}
\node[align=left](1) at (0,0) {
\begin{varwidth}{0.9\linewidth}
\end{varwidth}
};
\node[align=left](2) at (0,-1.5) {
\begin{varwidth}{0.9\linewidth}
\end{varwidth}
};
\path[#1,thick] (1) edge (2);
\end{tikzpicture}
\end{flushleft}
\vspace{-#3}
}
}

%Green check
\newcommand{\greencheck}{}%
\DeclareRobustCommand{\greencheck}{%
  \tikz\fill[scale=0.5, color=pastelgreen]
  (0,.35) -- (.25,0) -- (1,.7) -- (.25,.15) -- cycle;


% This is how you draw a graph in text:
\usetikzlibrary{shapes.misc}
\newcommand*{\graphdraw}{{\scriptsize
		\tikz[scale=0.1]{
			%%% edges
			\draw[thin] (2,0) -- (0,0) -- (0,2) -- (2,2);
			%% vertices
			\draw[fill=blue,draw=blue] (0,0) circle (12pt);
			\draw[fill=red,draw=red] (0,2) circle (12pt);
			\draw[fill=darkgreen,draw=darkgreen] (2,2) circle (12pt);
			\draw[fill=blue,draw=blue] (2,0) circle (12pt);
		}
}}}


% -===============                                    ===============%
% -===============                                    ===============%
% -===============                                    ===============%
%                         THEOREM ENVIRONMENTS                       %
% -===============                                    ===============%
% -===============                                    ===============%
% -===============                                    ===============%

\newtheorem{theorem}{Theorem}
\newtheorem{assumption}{Assumption}
\newtheorem{axiom}{Axiom}
\newtheorem{case}{Case}
\newtheorem{claim}{Claim}
\newtheorem{conclusion}{Conclusion}
\newtheorem{condition}{Condition}
\newtheorem{conjecture}{Conjecture}
\newtheorem{fact}{Fact}
\newtheorem{corollary}{Corollary}
\newtheorem{criterion}{Criterion}
\newtheorem{definition}{Definition}
\newtheorem{exercise}{Exercise}
\newtheorem{example}{Example}
\newtheorem{lemma}{Lemma}
\newtheorem{notation}{Notation}
\newtheorem{problem}{Problem}
\newtheorem{proposition}{Proposition}
\newtheorem{property}{Property}
\newtheorem{remark}{Remark}
\newtheorem{solution}{Solution}
\newtheorem{summary}{Summary}
\newtheorem{motivation}{Motivation}
\newtheorem{query}{Query}
\newtheorem{observation}{Observation}
% Let cleveref and thmtools work together
\makeatletter
\def\thmt@refnamewithcomma #1#2#3,#4,#5\@nil{%
	\@xa\def\csname\thmt@envname #1utorefname\endcsname{#3}%
	\ifcsname #2refname\endcsname
	\csname #2refname\expandafter\endcsname\expandafter{\thmt@envname}{#3}{#4}%
	\fi}
\makeatother
\Crefname{conjecture}{Conjecture}{Conjectures}
\Crefname{fact}{Fact}{Facts}
\Crefname{definition}{Definition}{Definitions}
\Crefname{observation}{Observation}{Observations}
\Crefname{assumption}{Assumption}{Assumptions}
\Crefname{axiom}{Axiom}{Axioms}
\Crefname{case}{Case}{Cases}
\Crefname{claim}{Claim}{Claims}
\Crefname{conclusion}{Conclusion}{Conclusions}
\Crefname{condition}{Condition}{Conditions}
\Crefname{criterion}{Criterion}{Criteria}
\Crefname{exercise}{Exercise}{Exercises}
\Crefname{example}{Example}{Examples}
\Crefname{notation}{Notation}{Notations}
\Crefname{problem}{Problem}{Problems}
\Crefname{property}{Property}{Properties}
\Crefname{remark}{Remark}{Remarks}
\Crefname{solution}{Solution}{Solutions}
\Crefname{summary}{Summary}{Summaries}
\Crefname{motivation}{Motivation}{Motivations}
\Crefname{query}{Query}{Queries}

% Styled boxes for prompts and callouts
\newtcblisting[auto counter, list inside=examplelist,Crefname={Prompt}{Prompts}]{prompt}[3][]{%
        colback=peach!50, colbacktitle=deepchampagne, coltitle=black,
        arc=0pt, titlerule=0, toptitle=2pt, bottomtitle=2pt,
        fonttitle=\bfseries,
        title={\small Prompt~\thetcbcounter:~#3},
        boxrule=0pt, opacityframe=0,
        listing only,
        label={#2},
        fontupper=\small,
        listing options={basicstyle=\ttfamily\small, breaklines=true},
        #1
}

\newmdenv[
  linecolor=black,
  linewidth=0.pt,
  roundcorner=10pt,
  backgroundcolor=customred!30,
]{attentionbox}

% -===============                                    ===============%
% -===============                                    ===============%
% -===============                                    ===============%
%                           MATH MACROS                              %
% -===============                                    ===============%
% -===============                                    ===============%
% -===============                                    ===============%
% general math ops
\newcommand{\mathcolorbox}[2]{\colorbox{#1}{$\displaystyle #2$}}
\newcommand{\one}{\mathds{1}} % indicator
\newcommand{\negquad}{\mkern-18mu}
\newcommand{\red}[1]{{\color{red} #1}}
\DeclareMathOperator*{\argmin}{argmin}
\DeclareMathOperator*{\argmax}{argmax}
\newcommand*\dbar[1]{\overline{\overline{\lower0.2ex\hbox{$#1$}}}}
\newcommand{\supp}{\text{supp\xspace}}
\DeclareMathOperator{\Tr}{Tr}
\newcommand{\dom}{\text{dom}}
\newcommand{\ran}{\text{ran}}
\newcommand{\setsize}[1]{\mid \! #1 \! \mid}
\newcommand{\norm}[1]{\mid\!\mid \! #1 \! \mid\!\mid}
\def\th{^{\text{th}}}
% stats — distributions, divergences, independence
\newcommand{\KL}{D_{\mathrm{KL}}}
\newcommand{\pr}{P}
\newcommand{\Exp}{\mathds{E}}
\newcommand{\eqdist}{\stackrel{d}{=}}
\newcommand{\eqas}{\stackrel{\text{a.s.}}{=}}
\newcommand{\standarderror}{\mathrm{SE}}
\newcommand{\indep}{\perp\!\!\!\perp}
\newcommand{\dtv}[0]{d_{\mathrm{TV}}}
\newcommand{\dkl}[0]{d_{\mathrm{KL}}}
\newcommand\bdot[1][.5]{\mathbin{\vcenter{\hbox{\scalebox{#1}{$\bullet$}}}}}
% do-calculus essentials
\newcommand{\doop}{\text{do}}
\newcommand{\Parents}{\text{Pa}}
\newcommand{\parents}{\text{pa}}
% graph theory
\newcommand{\iso}{\cong}
\newcommand{\notiso}{\not cong}
\newcommand{\aut}[1]{\text{Aut}(#1)}
% ML shorthands — saves typing in equations
\newcommand{\lr}{\alpha}
\newcommand{\reg}{\lambda}
\newcommand{\lf}{\mathcal{L}}
\newcommand{\sigmoid}{\sigma}
\newcommand{\softplus}{\zeta}
\newcommand{\relu}{\text{ReLU}}
% calligraphic alphabet — \cA through \cZ
\def\cA{{\mathcal{A}}}
\def\cB{{\mathcal{B}}}
\def\cC{{\mathcal{C}}}
\def\cD{{\mathcal{D}}}
\def\cE{{\mathcal{E}}}
\def\cF{{\mathcal{F}}}
\def\cG{{\mathcal{G}}}
\def\cH{{\mathcal{H}}}
\def\cI{{\mathcal{I}}}
\def\cJ{{\mathcal{J}}}
\def\cK{{\mathcal{K}}}
\def\cL{{\mathcal{L}}}
\def\cM{{\mathcal{M}}}
\def\cN{{\mathcal{N}}}
\def\cO{{\mathcal{O}}}
\def\cP{{\mathcal{P}}}
\def\cQ{{\mathcal{Q}}}
\def\cR{{\mathcal{R}}}
\def\cS{{\mathcal{S}}}
\def\cT{{\mathcal{T}}}
\def\cU{{\mathcal{U}}}
\def\cV{{\mathcal{V}}}
\def\cW{{\mathcal{W}}}
\def\cX{{\mathcal{X}}}
\def\cY{{\mathcal{Y}}}
\def\cZ{{\mathcal{Z}}}
% bold alphabet — \bA through \bZ, \ba through \bz
\def\bA{{\textbf{A}}}
\def\bB{{\textbf{B}}}
\def\bC{{\textbf{C}}}
\def\bD{{\textbf{D}}}
\def\bE{{\textbf{E}}}
\def\bF{{\textbf{F}}}
\def\bG{{\textbf{G}}}
\def\bH{{\textbf{H}}}
\def\bI{{\textbf{I}}}
\def\bJ{{\textbf{J}}}
\def\bK{{\textbf{K}}}
\def\bL{{\textbf{L}}}
\def\bM{{\textbf{M}}}
\def\bN{{\textbf{N}}}
\def\bO{{\textbf{O}}}
\def\bP{{\textbf{P}}}
\def\bQ{{\textbf{Q}}}
\def\bR{{\textbf{R}}}
\def\bS{{\textbf{S}}}
\def\bT{{\textbf{T}}}
\def\bU{{\textbf{U}}}
\def\bV{{\textbf{V}}}
\def\bW{{\textbf{W}}}
\def\bX{{\textbf{X}}}
\def\bY{{\textbf{Y}}}
\def\bZ{{\textbf{Z}}}
\def\ba{{\textbf{a}}}
\def\bb{{\textbf{b}}}
\def\bc{{\textbf{c}}}
\def\bd{{\textbf{d}}}
\def\be{{\textbf{e}}}
\def\bbf{{\textbf{f}}}
\def\bg{{\textbf{g}}}
\def\bh{{\textbf{h}}}
\def\bi{{\textbf{i}}}
\def\bj{{\textbf{j}}}
\def\bk{{\textbf{k}}}
\def\bl{{\textbf{l}}}
\def\bbm{{\textbf{m}}}
\def\bn{{\textbf{n}}}
\def\bo{{\textbf{o}}}
\def\bp{{\textbf{p}}}
\def\bq{{\textbf{q}}}
\def\br{{\textbf{r}}}
\def\bs{{\textbf{s}}}
\def\bt{{\textbf{t}}}
\def\bu{{\textbf{u}}}
\def\bv{{\textbf{v}}}
\def\bw{{\textbf{w}}}
\def\bx{{\textbf{x}}}
\def\by{{\textbf{y}}}
\def\bz{{\textbf{z}}}
% multiset braces: \lms x \rms → {{x}}
\newcommand*{\lms}{\{\mskip-5mu\{}
\newcommand*{\rms}{\}\mskip-5mu\}}


% -===============                                    ===============%
% -===============                                    ===============%
% -===============                                    ===============%
%                         MATH FONT COMMANDS                         %
% -===============                                    ===============%
% -===============                                    ===============%
% -===============                                    ===============%

% BOONDOX calligraphic — nicer than default \mathcal
\DeclareFontFamily{U}{BOONDOX-calo}{\skewchar\font=45 }
\DeclareFontShape{U}{BOONDOX-calo}{m}{n}{
  <-> s*[1.05] BOONDOX-r-calo}{}
\DeclareFontShape{U}{BOONDOX-calo}{b}{n}{
  <-> s*[1.05] BOONDOX-b-calo}{}
\DeclareMathAlphabet{\mathcalb}{U}{BOONDOX-calo}{m}{n}
\SetMathAlphabet{\mathcalb}{bold}{U}{BOONDOX-calo}{b}{n}
\DeclareMathAlphabet{\mathbcalb}{U}{BOONDOX-calo}{b}{n}
\newcommand{\sff}[1]{ \text{ {\sffamily \textbf{  {#1} }} } }
\newcommand{\ecal}[1]{ \EuScript{#1} }

% hyperref chokes on these in PDF metadata
\pdfstringdefDisableCommands{%
  \def\\{}%
  \def\texttt#1{<#1>}%
}

